\section{引言}

\subsection{平移扫描、窗口观测与诱导}

设 $G$ 为紧阿贝尔群，$m$ 为 Haar 测度，取 $a\in G$。定义平移
\[
T_a(x)=x+a,\qquad x\in G.
\]
给定可测窗口 $E\subset G$ 且 $m(E)>0$，定义首返时间
\[
r_E(x)=\min\{n\ge 1:\ T_a^n(x)\in E\},
\]
并定义诱导映射
\[
T_{a,E}(x)=T_a^{r_E(x)}(x),\qquad x\in E,
\]
以及归一化条件测度 $m_E(A)=m(A\cap E)/m(E)$。诱导系统 $(E,m_E,T_{a,E})$ 提取了“仅在窗口内可见”的动力学，并通过返回时间谱与分段结构体现算术与几何约束。

\subsection{研究问题}

本文围绕如下结构问题展开：
\begin{enumerate}
  \item 在圆周无理旋转上，区间窗口的返回时间谱如何由连分数算术决定，何时诱导仍保持为旋转类？
  \item 在二值稳定窗口族上，诱导—缩放的参数更新如何形成整数矩阵余循环，并与连分数尾部更新一致？二次无理数与黄金参数如何对应周期/不动点？
  \item 在高维环面平移中，平移向量的代数约束如何限制轨道闭包的有效维数，从而限制“单步扫描”可覆盖的维度？
  \item 如何以因子系统与范畴论语言将“观测—诱导—可观测压缩”组织为反射子与诱导函子之间的交换结构，并提出可检验推广命题？
\end{enumerate}

\subsection{贡献与改进点}

\begin{itemize}
  \item 将区间首返时间三值结构、二值退化判别、诱导 $2$-IET 等价旋转、诱导旋转数读出与分式线性表达组织为闭合定理链，得到可执行的“返回谱—诱导参数”算法化表达。
  \item 引入“窗口诱导重整化算子”，在二值稳定窗口族上给出参数更新的连分数尾部公式与矩阵作用表示，并据此刻画最终周期连分数参数的周期重整化轨道；黄金参数给出不动点与自诱导闭包链。
  \item 给出高维环面平移的闭包维数精确公式，并纠正“单步扫描在同一低次数数域内可填满高维”的误解：若坐标全部属于固定二次域，则闭包维数至多为 $1$。
  \item 提出反射子—诱导交换定理模板：在窗口可拉回与纤维随机性条件下，将诱导后的 Kronecker 因子与 Kronecker 因子上的诱导进行对齐，并将“可诱导闭包窗口”作为反射类内稳定对象形式化。
\end{itemize}

